```latex
\documentclass[11pt,a4paper]{article}

%====================================================
% ENCODAGE, LANGUE, MISE EN PAGE
%====================================================
\usepackage[utf8]{inputenc}
\usepackage[T1]{fontenc}
\usepackage[french]{babel}
\usepackage[a4paper,margin=2cm]{geometry}
\usepackage{lmodern}
\usepackage{microtype}
\usepackage{setspace}
\usepackage{fancyhdr}
\usepackage{lastpage}
\usepackage{titlesec}
\usepackage{enumitem}
\usepackage{array}
\usepackage{tabularx}
\usepackage{booktabs}
\usepackage{multicol}
\usepackage{hyperref}

%====================================================
% MATHS, PHYSIQUE, CHIMIE
%====================================================
\usepackage{amsmath,amssymb,amsthm}
\usepackage{siunitx}
\usepackage[version=4]{mhchem}
\usepackage{chemfig}
\sisetup{
  locale=FR,
  detect-all,
  per-mode=symbol,
  separate-uncertainty=true
}

%====================================================
% FIGURES ET ENCADRÉS
%====================================================
\usepackage{tikz}
\usetikzlibrary{arrows.meta,calc,positioning,patterns,decorations.pathmorphing}
\usepackage[most]{tcolorbox}

%====================================================
% COULEURS
%====================================================
\definecolor{bleu}{RGB}{25,70,130}
\definecolor{bleuclair}{RGB}{235,243,255}
\definecolor{vert}{RGB}{35,120,70}
\definecolor{vertclair}{RGB}{235,250,240}
\definecolor{orange}{RGB}{190,100,20}
\definecolor{orangeclair}{RGB}{255,246,232}
\definecolor{rouge}{RGB}{170,40,40}
\definecolor{rougeclair}{RGB}{255,238,238}
\definecolor{violet}{RGB}{95,55,150}
\definecolor{violetclair}{RGB}{246,240,255}
\definecolor{grisclair}{RGB}{246,246,246}

%====================================================
% EN-TÊTE
%====================================================
\pagestyle{fancy}
\fancyhf{}
\fancyhead[L]{Première spécialité physique-chimie}
\fancyhead[C]{TD -- Constitution et transformations de la matière}
\fancyhead[R]{Page \thepage/\pageref{LastPage}}
\fancyfoot[C]{\small Exercices progressifs avec corrigés détaillés}
\setlength{\headheight}{14pt}

%====================================================
% TITRES
%====================================================
\titleformat{\section}{\Large\bfseries\color{bleu}}{\thesection.}{0.5em}{}
\titleformat{\subsection}{\large\bfseries\color{violet}}{\thesubsection.}{0.5em}{}
\titleformat{\subsubsection}{\bfseries\color{orange}}{\thesubsubsection.}{0.5em}{}

%====================================================
% ENCADRÉS
%====================================================
\tcbset{
  enhanced,
  breakable,
  arc=2mm,
  boxrule=0.7pt,
  before skip=8pt,
  after skip=8pt
}

\newtcolorbox{methodbox}{
  colback=bleuclair,
  colframe=bleu,
  title=\textbf{Méthode},
  fonttitle=\bfseries
}

\newtcolorbox{retenirbox}{
  colback=vertclair,
  colframe=vert,
  title=\textbf{À retenir},
  fonttitle=\bfseries
}

\newtcolorbox{erreurbox}{
  colback=rougeclair,
  colframe=rouge,
  title=\textbf{Erreur classique},
  fonttitle=\bfseries
}

\newtcolorbox{vigilancebox}{
  colback=orangeclair,
  colframe=orange,
  title=\textbf{Point de vigilance},
  fonttitle=\bfseries
}

\newtcolorbox{correctionbox}{
  colback=grisclair,
  colframe=black!60,
  title=\textbf{Corrigé détaillé},
  fonttitle=\bfseries
}

%====================================================
% COMMANDES
%====================================================
\newcounter{exercice}

\newcommand{\etoileA}{$\star\circ\circ\circ\circ$}
\newcommand{\etoileB}{$\star\star\circ\circ\circ$}
\newcommand{\etoileC}{$\star\star\star\circ\circ$}
\newcommand{\etoileD}{$\star\star\star\star\circ$}
\newcommand{\etoileE}{$\star\star\star\star\star$}

\newcommand{\exo}[3]{
  \refstepcounter{exercice}
  \begin{tcolorbox}[
    colback=white,
    colframe=bleu,
    title=\textbf{Exercice \theexercice{} -- #1 -- #2},
    breakable
  ]
  #3
  \end{tcolorbox}
}

\newcommand{\corr}[2]{
  \begin{correctionbox}
  \textbf{Exercice #1.} #2
  \end{correctionbox}
}

\newcommand{\R}{\mathrm{R}}
\newcommand{\aq}{\mathrm{(aq)}}
\newcommand{\sol}{\mathrm{(s)}}
\newcommand{\liq}{\mathrm{(l)}}
\newcommand{\gaz}{\mathrm{(g)}}
\newcommand{\e}{\mathrm{e^-}}
\newcommand{\molL}{\si{mol.L^{-1}}}

%====================================================
% DOCUMENT
%====================================================
\begin{document}

%====================================================
% PAGE DE TITRE
%====================================================
\begin{titlepage}
\begin{center}

\vspace*{0.8cm}

{\Huge\bfseries TD complet}\\[0.3cm]
{\Huge\bfseries Constitution et transformations de la matière}\\[0.5cm]
{\Large Première générale -- Spécialité physique-chimie}\\[1cm]

\begin{tikzpicture}[scale=1]
  \draw[thick,bleu] (0,0) circle (1);
  \draw[thick,bleu] (3,0) circle (1);
  \draw[thick,bleu] (6,0) circle (1);
  \node at (0,0) {\Large $\ce{H2O}$};
  \node at (3,0) {\Large $\ce{Ox/Red}$};
  \node at (6,0) {\Large $\ce{C_xH_yO_z}$};
  \draw[-{Latex},thick,orange] (1.1,0) -- (1.9,0);
  \draw[-{Latex},thick,orange] (4.1,0) -- (4.9,0);
\end{tikzpicture}

\vfill

\begin{tcolorbox}[colback=bleuclair,colframe=bleu,width=0.92\textwidth,title=\textbf{Objectif général}]
Ce TD propose un entraînement progressif et approfondi sur les transformations chimiques,
les grandeurs molaires, l'avancement, les titrages, l'oxydo-réduction, la structure des entités,
la polarité, la cohésion, les extractions, la chimie organique, les synthèses et les combustions.
Il est conçu pour une classe de première spécialité physique-chimie de très bon niveau.
\end{tcolorbox}

\vfill

\begin{tabular}{rl}
\textbf{Nombre d'exercices :} & 50 exercices progressifs.\\
\textbf{Difficulté :} & de $\star\circ\circ\circ\circ$ à $\star\star\star\star\star$.\\
\textbf{Corrections :} & détaillées et rédigées à la fin du document.\\
\textbf{Utilisation :} & TD, entraînement, devoir maison, préparation aux contrôles.\\
\end{tabular}

\vfill

{\large \today}

\end{center}
\end{titlepage}

\tableofcontents
\newpage

%====================================================
\section{Rappels de méthode}

\subsection{Composition d'un système chimique}

\begin{retenirbox}
Pour un corps pur de masse $m$ et de masse molaire $M$ :
\[
n=\frac{m}{M}.
\]
Pour un gaz de volume $V$ et de volume molaire $V_m$ :
\[
n=\frac{V}{V_m}.
\]
Pour un soluté en solution de concentration $C$ et de volume $V$ :
\[
n=C V,
\]
avec $V$ en litre si $C$ est en \si{mol.L^{-1}}.
\end{retenirbox}

\begin{erreurbox}
Dans la relation $n=CV$, si $C$ est en \si{mol.L^{-1}}, le volume doit être exprimé en litre, pas en millilitre.
\[
\SI{25.0}{mL}=\SI{2.50e-2}{L}.
\]
\end{erreurbox}

\subsection{Avancement et tableau d'avancement}

\begin{methodbox}
Pour une réaction :
\[
aA+bB \longrightarrow cC+dD,
\]
si l'avancement est $x$, les variations de quantités de matière sont :
\[
\Delta n(A)=-ax,\quad \Delta n(B)=-bx,\quad \Delta n(C)=+cx,\quad \Delta n(D)=+dx.
\]
L'avancement maximal est la plus petite valeur imposée par les réactifs :
\[
x_{\max}=\min\left(\frac{n_0(A)}{a},\frac{n_0(B)}{b}\right).
\]
\end{methodbox}

\subsection{Oxydo-réduction}

\begin{retenirbox}
Un oxydant capte des électrons. Un réducteur cède des électrons.

Pour un couple oxydant-réducteur noté $\ce{Ox/Red}$ :
\[
\ce{Ox + n e- <=> Red}.
\]
Une réaction d'oxydo-réduction est obtenue en combinant deux demi-équations électroniques de manière à éliminer les électrons.
\end{retenirbox}

\subsection{Titrage}

\begin{retenirbox}
À l'équivalence d'un titrage, les réactifs titré et titrant ont été introduits en proportions stœchiométriques. Pour :
\[
aA+bB\longrightarrow \text{produits},
\]
si $A$ est titré et $B$ titrant, alors à l'équivalence :
\[
\frac{n(A)}{a}=\frac{n_E(B)}{b}.
\]
\end{retenirbox}

\subsection{Structure, polarité, cohésion}

\begin{methodbox}
Pour décider si une molécule est polaire :
\begin{enumerate}
  \item établir ou exploiter son schéma de Lewis ;
  \item déterminer la géométrie autour de l'atome central ;
  \item identifier les liaisons polarisées grâce aux électronégativités ;
  \item regarder si les polarités de liaison se compensent par symétrie.
\end{enumerate}
\end{methodbox}

\subsection{Chimie organique et combustion}

\begin{retenirbox}
Pour un alcane $\ce{C_nH_{2n+2}}$, la combustion complète s'écrit :
\[
\ce{C_nH_{2n+2} + \frac{3n+1}{2} O2 -> n CO2 + (n+1) H2O}.
\]
Le rendement d'une synthèse est :
\[
\eta=\frac{n_{\text{exp}}}{n_{\text{max}}}\times 100.
\]
\end{retenirbox}

\newpage

%====================================================
\section{Exercices}

%====================================================
\subsection{Niveau 1 : applications directes}

\exo{\etoileA}{Calculer une quantité de matière à partir d'une masse}{
On dispose de \SI{5.85}{g} de chlorure de sodium $\ce{NaCl}$.

Données : $M(\ce{Na})=\SI{23.0}{g.mol^{-1}}$, $M(\ce{Cl})=\SI{35.5}{g.mol^{-1}}$.

\begin{enumerate}
  \item Calculer la masse molaire de $\ce{NaCl}$.
  \item Calculer la quantité de matière de $\ce{NaCl}$.
  \item Donner le résultat avec un nombre adapté de chiffres significatifs.
\end{enumerate}
}

\exo{\etoileA}{Quantité de matière d'un gaz}{
Un ballon contient \SI{2.40}{L} de dioxygène $\ce{O2}$ dans des conditions où le volume molaire vaut $V_m=\SI{24.0}{L.mol^{-1}}$.

\begin{enumerate}
  \item Rappeler la relation entre $n$, $V$ et $V_m$.
  \item Calculer la quantité de matière de dioxygène.
  \item Calculer le nombre de molécules de dioxygène contenues dans le ballon.
\end{enumerate}

Donnée : $N_A=\SI{6.02e23}{mol^{-1}}$.
}

\exo{\etoileA}{Concentration en quantité de matière}{
On dissout \SI{0.0200}{mol} de glucose dans une fiole jaugée de volume \SI{250.0}{mL}.

\begin{enumerate}
  \item Convertir le volume en litre.
  \item Calculer la concentration en quantité de matière.
  \item Interpréter le résultat.
\end{enumerate}
}

\exo{\etoileA}{Masse à peser pour préparer une solution}{
On souhaite préparer \SI{100.0}{mL} d'une solution de sulfate de cuivre $\ce{CuSO4}$ de concentration $C=\SI{0.100}{mol.L^{-1}}$.

Données : $M(\ce{CuSO4})=\SI{159.6}{g.mol^{-1}}$.

\begin{enumerate}
  \item Calculer la quantité de matière nécessaire.
  \item Calculer la masse de sulfate de cuivre à peser.
\end{enumerate}
}

\exo{\etoileA}{Identifier oxydant et réducteur}{
On considère le couple $\ce{Cu^{2+}/Cu}$.

\begin{enumerate}
  \item Identifier l'oxydant du couple.
  \item Identifier le réducteur du couple.
  \item Écrire la demi-équation électronique associée.
\end{enumerate}
}

\exo{\etoileA}{Lecture d'un tableau d'avancement simple}{
On modélise la transformation par :
\[
\ce{A + 2B -> C}.
\]
À l'état initial, on dispose de $n_0(A)=\SI{0.20}{mol}$ et $n_0(B)=\SI{0.60}{mol}$.

\begin{enumerate}
  \item Écrire les quantités de matière en fonction de l'avancement $x$.
  \item Déterminer l'avancement maximal.
  \item Identifier le réactif limitant.
\end{enumerate}
}

\exo{\etoileA}{Équation de dissolution}{
Le chlorure de sodium solide se dissout dans l'eau.

\begin{enumerate}
  \item Écrire l'équation de dissolution de $\ce{NaCl(s)}$.
  \item Si l'on dissout \SI{0.10}{mol} de $\ce{NaCl}$ dans \SI{1.00}{L}, donner les concentrations en ions $\ce{Na+}$ et $\ce{Cl-}$.
\end{enumerate}
}

\exo{\etoileA}{Identifier un groupe caractéristique}{
Pour chacune des espèces suivantes, identifier la famille chimique :
\[
\ce{CH3CH2OH},\qquad \ce{CH3COOH},\qquad \ce{CH3CHO},\qquad \ce{CH3COCH3}.
\]

Familles possibles : alcool, acide carboxylique, aldéhyde, cétone.
}

\exo{\etoileA}{Rendement d'une synthèse}{
Une synthèse devrait produire au maximum \SI{0.120}{mol} d'un composé. On obtient expérimentalement \SI{0.096}{mol}.

\begin{enumerate}
  \item Rappeler la formule du rendement.
  \item Calculer le rendement.
  \item Interpréter le résultat.
\end{enumerate}
}

\exo{\etoileA}{Combustion complète du méthane}{
On étudie la combustion complète du méthane $\ce{CH4}$.

\begin{enumerate}
  \item Écrire l'équation ajustée de la combustion complète du méthane.
  \item Identifier les réactifs.
  \item Identifier les produits.
\end{enumerate}
}

%====================================================
\subsection{Niveau 2 : exercices standards}

\exo{\etoileB}{Mélange solide : composition molaire}{
Un alliage contient \SI{12.7}{g} de cuivre et \SI{6.50}{g} de zinc.

Données : $M(\ce{Cu})=\SI{63.5}{g.mol^{-1}}$, $M(\ce{Zn})=\SI{65.4}{g.mol^{-1}}$.

\begin{enumerate}
  \item Calculer la quantité de matière de cuivre.
  \item Calculer la quantité de matière de zinc.
  \item Déterminer la fraction molaire du cuivre.
\end{enumerate}
}

\exo{\etoileB}{Concentration massique et concentration molaire}{
Une solution de saccharose $\ce{C12H22O11}$ a une concentration en masse $c_m=\SI{34.2}{g.L^{-1}}$.

Données : $M(\ce{C})=\SI{12.0}{g.mol^{-1}}$, $M(\ce{H})=\SI{1.0}{g.mol^{-1}}$, $M(\ce{O})=\SI{16.0}{g.mol^{-1}}$.

\begin{enumerate}
  \item Calculer la masse molaire du saccharose.
  \item Déterminer la concentration en quantité de matière.
  \item Calculer la quantité de matière contenue dans \SI{200.0}{mL}.
\end{enumerate}
}

\exo{\etoileB}{Loi de Beer-Lambert}{
On mesure l'absorbance $A$ de solutions étalons d'une espèce colorée.

\begin{center}
\begin{tabular}{c|cccc}
$C$ (\si{mmol.L^{-1}}) & 0.50 & 1.00 & 1.50 & 2.00\\
\hline
$A$ & 0.12 & 0.24 & 0.36 & 0.48
\end{tabular}
\end{center}

Une solution inconnue présente une absorbance $A=0.30$.

\begin{enumerate}
  \item Justifier que les mesures suivent une loi de proportionnalité.
  \item Déterminer la concentration de la solution inconnue.
  \item Indiquer une limite expérimentale possible de la méthode.
\end{enumerate}
}

\exo{\etoileB}{Établir une réaction d'oxydo-réduction}{
On donne les couples :
\[
\ce{Fe^{2+}/Fe},\qquad \ce{Cu^{2+}/Cu}.
\]
On plonge une lame de fer dans une solution contenant des ions cuivre(II). La transformation observée est :
\[
\ce{Fe + Cu^{2+} -> Fe^{2+} + Cu}.
\]

\begin{enumerate}
  \item Écrire la demi-équation d'oxydation du fer.
  \item Écrire la demi-équation de réduction des ions cuivre(II).
  \item Vérifier l'équation globale.
  \item Identifier l'oxydant et le réducteur.
\end{enumerate}
}

\exo{\etoileB}{Avancement et réactif limitant}{
On fait réagir \SI{0.30}{mol} de dihydrogène avec \SI{0.20}{mol} de dioxygène selon :
\[
\ce{2H2 + O2 -> 2H2O}.
\]

\begin{enumerate}
  \item Construire le tableau d'avancement.
  \item Déterminer l'avancement maximal.
  \item Identifier le réactif limitant.
  \item Déterminer la composition finale.
\end{enumerate}
}

\exo{\etoileB}{Mélange stœchiométrique}{
On considère la réaction :
\[
\ce{2Al + 3S -> Al2S3}.
\]

\begin{enumerate}
  \item Donner la condition pour que le mélange initial soit stœchiométrique.
  \item On introduit \SI{0.40}{mol} d'aluminium. Quelle quantité de soufre faut-il introduire pour un mélange stœchiométrique ?
  \item Quelle quantité maximale de $\ce{Al2S3}$ peut alors se former ?
\end{enumerate}
}

\exo{\etoileB}{Titrage acido-basique simplifié}{
On titre \SI{20.0}{mL} d'une solution d'acide chlorhydrique $\ce{H3O+ + Cl-}$ par une solution d'hydroxyde de sodium de concentration $C_B=\SI{0.100}{mol.L^{-1}}$.
L'équivalence est obtenue pour $V_E=\SI{12.5}{mL}$.

La réaction support est :
\[
\ce{H3O+ + HO- -> 2H2O}.
\]

\begin{enumerate}
  \item Écrire la relation à l'équivalence.
  \item Calculer la quantité de matière d'ions hydroxyde introduits à l'équivalence.
  \item En déduire la concentration de la solution acide.
\end{enumerate}
}

\exo{\etoileB}{Lewis et géométrie de l'eau}{
On étudie la molécule d'eau $\ce{H2O}$.

\begin{enumerate}
  \item Donner le nombre d'électrons de valence de l'oxygène et de l'hydrogène.
  \item Établir le schéma de Lewis de $\ce{H2O}$.
  \item Préciser la géométrie de la molécule.
  \item La molécule est-elle polaire ? Justifier.
\end{enumerate}
}

\exo{\etoileB}{Dissolution du chlorure de calcium}{
Le chlorure de calcium solide se dissout selon :
\[
\ce{CaCl2(s) -> Ca^{2+}(aq) + 2Cl-(aq)}.
\]
On dissout \SI{0.050}{mol} de $\ce{CaCl2}$ dans de l'eau pour obtenir \SI{500.0}{mL} de solution.

\begin{enumerate}
  \item Calculer la concentration en soluté apporté.
  \item Calculer la concentration en ions calcium.
  \item Calculer la concentration en ions chlorure.
\end{enumerate}
}

\exo{\etoileB}{Nom et formule d'un alcool}{
On considère le propan-1-ol, de formule semi-développée :
\[
\ce{CH3-CH2-CH2-OH}.
\]

\begin{enumerate}
  \item Identifier la chaîne carbonée principale.
  \item Identifier le groupe caractéristique.
  \item Justifier le nom propan-1-ol.
\end{enumerate}
}

%====================================================
\subsection{Niveau 3 : exercices intermédiaires}

\exo{\etoileC}{Préparation d'une gamme étalon}{
On souhaite préparer \SI{100.0}{mL} de solutions étalons de permanganate de potassium de concentrations :
\[
\SI{1.0e-4}{mol.L^{-1}},\quad
\SI{2.0e-4}{mol.L^{-1}},\quad
\SI{4.0e-4}{mol.L^{-1}}.
\]
On dispose d'une solution mère de concentration $C_0=\SI{1.0e-3}{mol.L^{-1}}$.

\begin{enumerate}
  \item Rappeler la relation de dilution.
  \item Calculer le volume de solution mère à prélever pour chaque solution fille.
  \item Proposer un protocole expérimental pour préparer la solution à $\SI{2.0e-4}{mol.L^{-1}}$.
\end{enumerate}
}

\exo{\etoileC}{Beer-Lambert et choix de longueur d'onde}{
Une espèce colorée absorbe fortement à \SI{520}{nm}. Son spectre montre une absorbance maximale dans le vert.

\begin{enumerate}
  \item Quelle couleur est principalement absorbée ?
  \item Quelle couleur la solution peut-elle transmettre ou diffuser majoritairement ?
  \item Pourquoi choisit-on souvent la longueur d'onde du maximum d'absorption pour réaliser un dosage ?
  \item On mesure $A=0.65$ pour une solution inconnue. Sachant que $A=1250C$ avec $C$ en \si{mol.L^{-1}}, déterminer $C$.
\end{enumerate}
}

\exo{\etoileC}{Détartrage par un acide}{
Le carbonate de calcium $\ce{CaCO3}$ réagit avec les ions oxonium selon :
\[
\ce{CaCO3(s) + 2H3O+(aq) -> Ca^{2+}(aq) + CO2(g) + 3H2O(l)}.
\]
On introduit \SI{0.010}{mol} de $\ce{CaCO3}$ et \SI{0.015}{mol} de $\ce{H3O+}$.

\begin{enumerate}
  \item Construire le tableau d'avancement.
  \item Déterminer l'avancement maximal.
  \item Identifier le réactif limitant.
  \item Déterminer la composition finale.
\end{enumerate}
}

\exo{\etoileC}{Transformation non totale}{
La transformation :
\[
\ce{A + B -> C}
\]
est réalisée avec $n_0(A)=\SI{0.50}{mol}$ et $n_0(B)=\SI{0.40}{mol}$. À l'état final, il reste \SI{0.20}{mol} de $B$.

\begin{enumerate}
  \item Déterminer l'avancement final $x_f$.
  \item Déterminer l'avancement maximal $x_{\max}$.
  \item La transformation est-elle totale ?
  \item Déterminer les quantités finales de $A$ et $C$.
\end{enumerate}
}

\exo{\etoileC}{Équilibrer une réaction redox en milieu acide}{
On donne les couples :
\[
\ce{MnO4-/Mn^{2+}},\qquad \ce{Fe^{3+}/Fe^{2+}}.
\]
En milieu acide, les demi-équations sont :
\[
\ce{MnO4- + 8H+ + 5e- -> Mn^{2+} + 4H2O},
\]
\[
\ce{Fe^{2+} -> Fe^{3+} + e-}.
\]

\begin{enumerate}
  \item Multiplier les demi-équations pour éliminer les électrons.
  \item Écrire l'équation globale.
  \item Identifier l'oxydant et le réducteur.
\end{enumerate}
}

\exo{\etoileC}{Titrage du fer(II) par le permanganate}{
On titre \SI{10.0}{mL} d'une solution contenant des ions $\ce{Fe^{2+}}$ par une solution de permanganate de concentration $C_P=\SI{2.00e-2}{mol.L^{-1}}$.
L'équivalence est obtenue pour $V_E=\SI{14.0}{mL}$.

La réaction est :
\[
\ce{MnO4- + 5Fe^{2+} + 8H+ -> Mn^{2+} + 5Fe^{3+} + 4H2O}.
\]

\begin{enumerate}
  \item Écrire la relation stœchiométrique à l'équivalence.
  \item Calculer la quantité de permanganate introduite à l'équivalence.
  \item Déterminer la quantité initiale de $\ce{Fe^{2+}}$.
  \item Calculer la concentration en $\ce{Fe^{2+}}$.
\end{enumerate}
}

\exo{\etoileC}{Géométrie et polarité du dioxyde de carbone}{
On étudie la molécule $\ce{CO2}$.

\begin{enumerate}
  \item Établir son schéma de Lewis.
  \item Préciser sa géométrie.
  \item Les liaisons C=O sont-elles polarisées ?
  \item La molécule $\ce{CO2}$ est-elle polaire ? Justifier soigneusement.
\end{enumerate}
}

\exo{\etoileC}{Polarité de l'ammoniac}{
On étudie la molécule $\ce{NH3}$.

\begin{enumerate}
  \item Établir le schéma de Lewis de $\ce{NH3}$.
  \item Préciser la géométrie autour de l'atome d'azote.
  \item Expliquer pourquoi la molécule est polaire.
  \item Comparer qualitativement avec $\ce{BF3}$, molécule plane triangulaire apolaire.
\end{enumerate}
}

\exo{\etoileC}{Extraction de l'iode}{
Le diiode $\ce{I2}$ est peu soluble dans l'eau mais très soluble dans le cyclohexane, solvant organique non miscible à l'eau. On dispose d'une solution aqueuse contenant du diiode.

\begin{enumerate}
  \item Quel solvant peut-on utiliser pour extraire le diiode ?
  \item Expliquer le rôle de l'ampoule à décanter.
  \item Dans quelle phase se trouvera majoritairement le diiode après agitation et décantation ?
  \item Quel critère de solubilité est illustré ici ?
\end{enumerate}
}

\exo{\etoileC}{Combustion de l'éthanol}{
L'éthanol a pour formule $\ce{C2H6O}$.

\begin{enumerate}
  \item Écrire l'équation ajustée de sa combustion complète.
  \item Calculer la quantité de dioxyde de carbone formée lors de la combustion de \SI{0.50}{mol} d'éthanol.
  \item Calculer la quantité de dioxygène consommée.
  \item Commenter l'impact environnemental du dioxyde de carbone formé.
\end{enumerate}
}

%====================================================
\subsection{Niveau 4 : exercices difficiles}

\exo{\etoileD}{Contrôle qualité d'un comprimé de vitamine C}{
Un comprimé de vitamine C contient de l'acide ascorbique $\ce{C6H8O6}$. On dissout un comprimé dans une fiole jaugée de \SI{250.0}{mL}. On dose ensuite \SI{20.0}{mL} de cette solution par une solution de diiode de concentration $C_I=\SI{5.00e-3}{mol.L^{-1}}$.
L'équivalence est obtenue pour $V_E=\SI{18.2}{mL}$.

La réaction support est supposée totale et de proportion :
\[
\ce{C6H8O6 + I2 -> C6H6O6 + 2I- + 2H+}.
\]

Donnée : $M(\ce{C6H8O6})=\SI{176.0}{g.mol^{-1}}$.

\begin{enumerate}
  \item Déterminer la quantité de diiode introduite à l'équivalence.
  \item En déduire la quantité d'acide ascorbique dans le prélèvement.
  \item En déduire la quantité d'acide ascorbique dans la fiole.
  \item Calculer la masse d'acide ascorbique contenue dans le comprimé.
  \item Comparer à une indication fabricant de \SI{500}{mg}.
\end{enumerate}
}

\exo{\etoileD}{Analyse d'une eau ferrugineuse}{
Une eau contient des ions $\ce{Fe^{2+}}$. On prélève \SI{50.0}{mL} d'eau et on titre les ions $\ce{Fe^{2+}}$ par une solution de permanganate de potassium acidifiée de concentration $C=\SI{1.00e-3}{mol.L^{-1}}$.
L'équivalence est obtenue pour $V_E=\SI{9.60}{mL}$.

La réaction est :
\[
\ce{MnO4- + 5Fe^{2+} + 8H+ -> Mn^{2+} + 5Fe^{3+} + 4H2O}.
\]

\begin{enumerate}
  \item Expliquer pourquoi le permanganate peut servir d'indicateur coloré.
  \item Calculer la quantité de $\ce{MnO4-}$ versée à l'équivalence.
  \item Calculer la concentration molaire des ions $\ce{Fe^{2+}}$ dans l'eau.
  \item Calculer la concentration massique en fer.
\end{enumerate}

Donnée : $M(\ce{Fe})=\SI{55.8}{g.mol^{-1}}$.
}

\exo{\etoileD}{Méthode numérique : état final d'une transformation totale}{
On considère la réaction totale :
\[
\ce{2A + 3B -> C + 4D}.
\]
On introduit $n_0(A)=\SI{0.80}{mol}$ et $n_0(B)=\SI{0.90}{mol}$.

\begin{enumerate}
  \item Écrire les expressions des quantités finales en fonction de $x$.
  \item Déterminer $x_{\max}$.
  \item En déduire la composition finale.
  \item Proposer un court algorithme en pseudo-code permettant de déterminer automatiquement le réactif limitant pour une réaction $aA+bB\to$ produits.
\end{enumerate}
}

\exo{\etoileD}{Spectre d'absorption et couleur perçue}{
Une solution présente le spectre d'absorption simplifié suivant :

\begin{center}
\begin{tabular}{c|cccccc}
$\lambda$ (\si{nm}) & 400 & 450 & 500 & 550 & 600 & 650\\
\hline
$A$ & 0.10 & 0.25 & 0.90 & 1.20 & 0.35 & 0.10
\end{tabular}
\end{center}

\begin{enumerate}
  \item Dans quel domaine de longueurs d'onde l'absorption est-elle maximale ?
  \item Quelle couleur est principalement absorbée ?
  \item Prévoir qualitativement la couleur de la solution.
  \item Pourquoi le choix de $\lambda=\SI{550}{nm}$ est-il judicieux pour un dosage par absorbance ?
\end{enumerate}
}

\exo{\etoileD}{Molécule avec lacune électronique : le borane}{
On étudie la molécule $\ce{BH3}$.

\begin{enumerate}
  \item Donner le nombre d'électrons de valence du bore et de l'hydrogène.
  \item Établir le schéma de Lewis de $\ce{BH3}$.
  \item Montrer que le bore ne respecte pas l'octet.
  \item Expliquer ce que signifie l'expression « lacune électronique ».
  \item La molécule, plane triangulaire, est-elle polaire ?
\end{enumerate}
}

\exo{\etoileD}{Dissolution d'un solide ionique et concentrations ioniques}{
On dissout \SI{3.42}{g} de sulfate d'aluminium $\ce{Al2(SO4)3}$ dans l'eau et on complète à \SI{250.0}{mL}.

Données :
\[
M(\ce{Al})=\SI{27.0}{g.mol^{-1}},\quad
M(\ce{S})=\SI{32.1}{g.mol^{-1}},\quad
M(\ce{O})=\SI{16.0}{g.mol^{-1}}.
\]

\begin{enumerate}
  \item Écrire l'équation de dissolution.
  \item Calculer la masse molaire de $\ce{Al2(SO4)3}$.
  \item Calculer la concentration en soluté apporté.
  \item Calculer les concentrations des ions en solution.
\end{enumerate}
}

\exo{\etoileD}{Savon et amphiphilie}{
On modélise un ion savon par une longue chaîne carbonée hydrophobe terminée par un groupe carboxylate hydrophile :
\[
\ce{CH3-(CH2)16-COO-}.
\]

\begin{enumerate}
  \item Identifier la partie hydrophobe.
  \item Identifier la partie hydrophile.
  \item Expliquer pourquoi cette entité est amphiphile.
  \item Expliquer qualitativement comment un savon peut aider à enlever une tache grasse.
\end{enumerate}
}

\exo{\etoileD}{Analyse d'un spectre infrarouge}{
Une molécule organique de formule brute $\ce{C3H6O}$ donne en spectroscopie infrarouge une forte bande vers \SI{1715}{cm^{-1}} et aucune large bande vers \SI{3200}{cm^{-1}}.

Données :
\begin{itemize}
  \item liaison O--H alcool : large bande vers \SI{3200}{cm^{-1}} à \SI{3600}{cm^{-1}} ;
  \item liaison C=O : forte bande vers \SI{1650}{cm^{-1}} à \SI{1750}{cm^{-1}}.
\end{itemize}

\begin{enumerate}
  \item Quelle liaison est mise en évidence ?
  \item Quelle famille peut-on exclure ?
  \item Proposer deux familles compatibles avec la formule et le spectre.
  \item Proposer une formule semi-développée possible.
\end{enumerate}
}

\exo{\etoileD}{Rendement d'une synthèse organique}{
On synthétise un ester à partir de \SI{0.20}{mol} d'acide carboxylique et \SI{0.50}{mol} d'alcool. La réaction est supposée de proportion $1:1$. Après purification, on obtient \SI{17.6}{g} d'ester de masse molaire $M=\SI{116}{g.mol^{-1}}$.

\begin{enumerate}
  \item Identifier le réactif limitant.
  \item Calculer la quantité maximale d'ester attendue.
  \item Calculer la quantité d'ester obtenue.
  \item Calculer le rendement.
\end{enumerate}
}

\exo{\etoileD}{Énergie de combustion à partir des liaisons}{
On estime l'énergie molaire de combustion du méthane gazeux :
\[
\ce{CH4(g) + 2O2(g) -> CO2(g) + 2H2O(g)}.
\]

Énergies de liaison moyennes :
\[
E(\ce{C-H})=\SI{413}{kJ.mol^{-1}},\quad
E(\ce{O=O})=\SI{498}{kJ.mol^{-1}},
\]
\[
E(\ce{C=O})=\SI{799}{kJ.mol^{-1}},\quad
E(\ce{O-H})=\SI{463}{kJ.mol^{-1}}.
\]

\begin{enumerate}
  \item Identifier les liaisons rompues.
  \item Identifier les liaisons formées.
  \item Estimer l'énergie molaire de réaction.
  \item La combustion est-elle exothermique ?
\end{enumerate}
}

%====================================================
\subsection{Niveau 5 : exercices ambitieux et problèmes longs}

\exo{\etoileE}{Problème long : détartrage d'une bouilloire}{
Une bouilloire contient un dépôt de tartre assimilé à du carbonate de calcium $\ce{CaCO3}$. On verse \SI{100.0}{mL} d'une solution d'acide chlorhydrique de concentration $C=\SI{1.00}{mol.L^{-1}}$.
La réaction est :
\[
\ce{CaCO3(s) + 2H3O+(aq) -> Ca^{2+}(aq) + CO2(g) + 3H2O(l)}.
\]
Après réaction, tout le tartre a disparu et il reste \SI{0.020}{mol} d'ions $\ce{H3O+}$.

Données : $M(\ce{CaCO3})=\SI{100.1}{g.mol^{-1}}$.

\begin{enumerate}
  \item Calculer la quantité initiale d'ions $\ce{H3O+}$.
  \item Déterminer l'avancement final.
  \item Déterminer la quantité initiale de carbonate de calcium.
  \item Calculer la masse de tartre éliminée.
  \item Calculer le volume de $\ce{CO2}$ formé si $V_m=\SI{24.0}{L.mol^{-1}}$.
  \item Expliquer pourquoi un dégagement gazeux est observé.
\end{enumerate}
}

\exo{\etoileE}{Problème long : dosage colorimétrique d'un colorant alimentaire}{
On veut déterminer la concentration d'un colorant bleu dans une boisson. On prépare une gamme étalon :

\begin{center}
\begin{tabular}{c|ccccc}
$C$ (\si{mg.L^{-1}}) & 2.0 & 4.0 & 6.0 & 8.0 & 10.0\\
\hline
$A$ & 0.15 & 0.31 & 0.45 & 0.60 & 0.76
\end{tabular}
\end{center}

La boisson diluée 5 fois présente une absorbance $A=0.54$.

\begin{enumerate}
  \item Justifier que la loi de Beer-Lambert peut être utilisée approximativement.
  \item Déterminer le coefficient directeur moyen de la droite $A=kC$.
  \item Déterminer la concentration du colorant dans la boisson diluée.
  \item En déduire la concentration dans la boisson initiale.
  \item Expliquer pourquoi il peut être nécessaire de diluer la boisson avant mesure.
\end{enumerate}
}

\exo{\etoileE}{Problème long : titrage d'une eau oxygénée}{
L'eau oxygénée contient du peroxyde d'hydrogène $\ce{H2O2}$. On titre \SI{10.0}{mL} d'une solution diluée d'eau oxygénée par une solution de permanganate de potassium acidifiée de concentration $C=\SI{2.00e-2}{mol.L^{-1}}$.
L'équivalence est obtenue pour $V_E=\SI{16.0}{mL}$.

La réaction support est :
\[
\ce{2MnO4- + 5H2O2 + 6H+ -> 2Mn^{2+} + 5O2 + 8H2O}.
\]

\begin{enumerate}
  \item Identifier l'oxydant titrant.
  \item Écrire la relation stœchiométrique à l'équivalence.
  \item Calculer la quantité de permanganate versée.
  \item En déduire la quantité de peroxyde d'hydrogène dans le prélèvement.
  \item Calculer la concentration molaire de $\ce{H2O2}$ dans la solution diluée.
  \item Si la solution commerciale a été diluée 10 fois, déterminer sa concentration commerciale.
\end{enumerate}
}

\exo{\etoileE}{Problème long : corrosion du fer}{
La corrosion du fer peut être modélisée par une étape d'oxydation :
\[
\ce{Fe -> Fe^{2+} + 2e-}.
\]
Le dioxygène dissous peut être réduit selon :
\[
\ce{O2 + 4H+ + 4e- -> 2H2O}.
\]

\begin{enumerate}
  \item Écrire l'équation globale en combinant les deux demi-équations.
  \item Identifier l'oxydant et le réducteur.
  \item Pour \SI{1.12}{g} de fer oxydé, calculer la quantité de fer consommée.
  \item Calculer la quantité de dioxygène nécessaire.
  \item Calculer le volume de dioxygène correspondant avec $V_m=\SI{24.0}{L.mol^{-1}}$.
  \item Expliquer pourquoi la corrosion est favorisée en milieu humide.
\end{enumerate}

Donnée : $M(\ce{Fe})=\SI{55.8}{g.mol^{-1}}$.
}

\exo{\etoileE}{Problème long : extraction d'une espèce odorante}{
Une espèce odorante $S$ est dissoute dans \SI{100}{mL} d'eau. Elle est peu soluble dans l'eau mais très soluble dans l'éthanoate d'éthyle, non miscible à l'eau.

On donne les solubilités approximatives :
\[
s_{\text{eau}}=\SI{0.50}{g.L^{-1}},\qquad
s_{\text{org}}=\SI{25}{g.L^{-1}}.
\]

\begin{enumerate}
  \item Quel solvant choisir pour extraire $S$ ?
  \item Expliquer pourquoi les deux liquides doivent être non miscibles.
  \item Après agitation et décantation, dans quelle phase se trouve majoritairement $S$ ?
  \item Pourquoi réalise-t-on souvent plusieurs extractions successives avec de petits volumes plutôt qu'une seule extraction ?
  \item Proposer les étapes expérimentales d'une extraction liquide-liquide.
\end{enumerate}
}

\exo{\etoileE}{Problème long : polarité, solubilité et miscibilité}{
On compare trois espèces : l'eau $\ce{H2O}$, l'éthanol $\ce{CH3CH2OH}$ et l'hexane $\ce{C6H14}$.

\begin{enumerate}
  \item Expliquer pourquoi l'eau est une molécule polaire.
  \item Identifier dans l'éthanol la partie polaire et la partie apolaire.
  \item Expliquer pourquoi l'éthanol est miscible à l'eau.
  \item Expliquer pourquoi l'hexane n'est pas miscible à l'eau.
  \item Prévoir dans quel solvant, eau ou hexane, une espèce très apolaire sera la plus soluble.
  \item Relier ces observations à l'expression « qui se ressemble s'assemble ».
\end{enumerate}
}

\exo{\etoileE}{Problème long : synthèse, purification et rendement}{
Lors de la synthèse d'un ester, on introduit \SI{0.150}{mol} d'acide carboxylique et \SI{0.300}{mol} d'alcool. La réaction est de proportion $1:1$.
Après chauffage à reflux, lavage, séchage et purification, on obtient \SI{12.8}{g} d'ester. La masse molaire de l'ester est $M=\SI{102}{g.mol^{-1}}$.

\begin{enumerate}
  \item Identifier les étapes de transformation, d'isolement, de purification et d'analyse.
  \item Déterminer le réactif limitant.
  \item Calculer la quantité maximale d'ester attendue.
  \item Calculer la quantité réelle obtenue.
  \item Calculer le rendement.
  \item Donner deux raisons expérimentales pouvant expliquer un rendement inférieur à 100\%.
\end{enumerate}
}

\exo{\etoileE}{Problème long : nomenclature et spectroscopie IR}{
On dispose d'une molécule organique à chaîne saturée comportant trois atomes de carbone et un seul groupe caractéristique. Son spectre IR présente une large bande intense entre \SI{3200}{cm^{-1}} et \SI{3600}{cm^{-1}}, mais pas de bande forte vers \SI{1700}{cm^{-1}}.

\begin{enumerate}
  \item Quelle liaison est mise en évidence par la bande large ?
  \item Quelle famille chimique est probable ?
  \item Proposer deux formules semi-développées compatibles.
  \item Donner les noms correspondants.
  \item Ces deux molécules ont-elles la même formule brute ?
  \item Comment appelle-t-on de telles espèces ?
\end{enumerate}
}

\exo{\etoileE}{Problème long : pouvoir calorifique d'un alcool}{
On brûle \SI{2.30}{g} d'éthanol pour chauffer \SI{500}{g} d'eau. La température de l'eau augmente de \SI{18.0}{\celsius}. On suppose que toute l'énergie libérée chauffe l'eau.

Données :
\[
c_{\text{eau}}=\SI{4.18}{J.g^{-1}.\celsius^{-1}},\qquad M(\ce{C2H6O})=\SI{46.0}{g.mol^{-1}}.
\]

\begin{enumerate}
  \item Calculer l'énergie reçue par l'eau.
  \item Calculer la quantité de matière d'éthanol brûlée.
  \item Estimer l'énergie libérée par mole d'éthanol.
  \item Estimer le pouvoir calorifique massique expérimental en \si{kJ.g^{-1}}.
  \item Pourquoi cette estimation est-elle souvent inférieure à la valeur tabulée ?
\end{enumerate}
}

\exo{\etoileE}{Problème concours : batterie, redox et bilan de matière}{
Une pile simplifiée met en jeu les couples $\ce{Zn^{2+}/Zn}$ et $\ce{Cu^{2+}/Cu}$. La réaction globale est :
\[
\ce{Zn(s) + Cu^{2+}(aq) -> Zn^{2+}(aq) + Cu(s)}.
\]
On introduit une lame de zinc de masse \SI{1.31}{g} dans \SI{100.0}{mL} d'une solution de sulfate de cuivre(II) de concentration $C=\SI{0.150}{mol.L^{-1}}$.

Données :
\[
M(\ce{Zn})=\SI{65.4}{g.mol^{-1}},\qquad M(\ce{Cu})=\SI{63.5}{g.mol^{-1}}.
\]

\begin{enumerate}
  \item Identifier l'oxydant et le réducteur.
  \item Écrire les demi-équations électroniques.
  \item Calculer les quantités initiales de zinc et d'ions cuivre(II).
  \item Déterminer le réactif limitant.
  \item Déterminer la masse de cuivre formée.
  \item Déterminer la masse de zinc restante.
  \item Expliquer qualitativement l'origine du courant électrique dans une pile.
\end{enumerate}
}

\newpage

%====================================================
\section{Exercices bilan}

\exo{\etoileD}{Bilan 1 -- Analyse complète d'une solution colorée}{
Une solution colorée contient des ions $\ce{Cu^{2+}}$. On prépare une gamme étalon et on obtient la relation :
\[
A=120C,
\]
où $C$ est en \si{mol.L^{-1}}. Une solution inconnue diluée 4 fois présente une absorbance $A=0.36$.

\begin{enumerate}
  \item Déterminer la concentration de la solution diluée.
  \item Déterminer la concentration de la solution initiale.
  \item Calculer la quantité d'ions $\ce{Cu^{2+}}$ dans \SI{250.0}{mL} de solution initiale.
  \item Si ces ions proviennent de la dissolution de $\ce{CuSO4}$, calculer la masse de sulfate de cuivre anhydre dissoute.
\end{enumerate}

Donnée : $M(\ce{CuSO4})=\SI{159.6}{g.mol^{-1}}$.
}

\exo{\etoileE}{Bilan 2 -- Titrage, avancement et rendement}{
On synthétise un acide organique solide. On en prélève \SI{0.620}{g}, que l'on dissout dans l'eau. On titre la solution obtenue par une solution d'hydroxyde de sodium de concentration $C_B=\SI{0.200}{mol.L^{-1}}$.
L'équivalence est atteinte pour $V_E=\SI{25.0}{mL}$. L'acide est monoprotique et réagit selon :
\[
\ce{AH + HO- -> A- + H2O}.
\]
La synthèse aurait dû donner au maximum \SI{0.00620}{mol} d'acide.

\begin{enumerate}
  \item Calculer la quantité de base versée à l'équivalence.
  \item En déduire la quantité d'acide dans l'échantillon.
  \item Déterminer la masse molaire expérimentale de l'acide.
  \item Calculer le rendement molaire de la synthèse.
  \item Proposer deux contrôles expérimentaux permettant de vérifier la pureté du produit.
\end{enumerate}
}

\newpage

%====================================================
\section{Corrigés détaillés}

\corr{1}{
\begin{enumerate}
  \item La masse molaire est :
  \[
  M(\ce{NaCl})=M(\ce{Na})+M(\ce{Cl})=23.0+35.5=\SI{58.5}{g.mol^{-1}}.
  \]
  \item On utilise :
  \[
  n=\frac{m}{M}.
  \]
  Donc :
  \[
  n=\frac{5.85}{58.5}=\SI{0.100}{mol}.
  \]
  \item Le résultat est cohérent avec trois chiffres significatifs : $n=\SI{0.100}{mol}$.
\end{enumerate}
}

\corr{2}{
\begin{enumerate}
  \item Pour un gaz :
  \[
  n=\frac{V}{V_m}.
  \]
  \item
  \[
  n=\frac{2.40}{24.0}=\SI{0.100}{mol}.
  \]
  \item Le nombre de molécules vaut :
  \[
  N=nN_A=0.100\times 6.02\times 10^{23}=\SI{6.02e22}{}.
  \]
\end{enumerate}
}

\corr{3}{
\begin{enumerate}
  \item
  \[
  V=\SI{250.0}{mL}=\SI{0.2500}{L}.
  \]
  \item
  \[
  C=\frac{n}{V}=\frac{0.0200}{0.2500}=\SI{0.0800}{mol.L^{-1}}.
  \]
  \item La solution contient \SI{0.0800}{mol} de glucose par litre de solution.
\end{enumerate}
}

\corr{4}{
\begin{enumerate}
  \item
  \[
  V=\SI{100.0}{mL}=\SI{0.1000}{L}.
  \]
  \[
  n=CV=0.100\times 0.1000=\SI{1.00e-2}{mol}.
  \]
  \item
  \[
  m=nM=1.00\times 10^{-2}\times 159.6=\SI{1.60}{g}.
  \]
\end{enumerate}
}

\corr{5}{
\begin{enumerate}
  \item Dans le couple $\ce{Cu^{2+}/Cu}$, l'oxydant est $\ce{Cu^{2+}}$.
  \item Le réducteur est $\ce{Cu}$.
  \item La demi-équation est :
  \[
  \ce{Cu^{2+} + 2e- -> Cu}.
  \]
\end{enumerate}
}

\corr{6}{
\begin{enumerate}
  \item Pour $\ce{A + 2B -> C}$ :
  \[
  n(A)=0.20-x,\quad n(B)=0.60-2x,\quad n(C)=x.
  \]
  \item Les contraintes sont :
  \[
  0.20-x\geq 0 \Rightarrow x\leq 0.20,
  \]
  \[
  0.60-2x\geq 0 \Rightarrow x\leq 0.30.
  \]
  Donc :
  \[
  x_{\max}=\SI{0.20}{mol}.
  \]
  \item Le réactif limitant est $A$.
\end{enumerate}
}

\corr{7}{
\begin{enumerate}
  \item
  \[
  \ce{NaCl(s) -> Na+(aq) + Cl-(aq)}.
  \]
  \item Dans \SI{1.00}{L}, si $n=\SI{0.10}{mol}$ :
  \[
  [\ce{Na+}]=[\ce{Cl-}]=\SI{0.10}{mol.L^{-1}}.
  \]
\end{enumerate}
}

\corr{8}{
\[
\ce{CH3CH2OH} : \text{alcool}.
\]
\[
\ce{CH3COOH} : \text{acide carboxylique}.
\]
\[
\ce{CH3CHO} : \text{aldéhyde}.
\]
\[
\ce{CH3COCH3} : \text{cétone}.
\]
}

\corr{9}{
\begin{enumerate}
  \item
  \[
  \eta=\frac{n_{\text{obtenu}}}{n_{\max}}\times 100.
  \]
  \item
  \[
  \eta=\frac{0.096}{0.120}\times 100=80.0\%.
  \]
  \item On a obtenu 80\% de la quantité maximale théorique.
\end{enumerate}
}

\corr{10}{
\begin{enumerate}
  \item
  \[
  \ce{CH4 + 2O2 -> CO2 + 2H2O}.
  \]
  \item Les réactifs sont le méthane et le dioxygène.
  \item Les produits sont le dioxyde de carbone et l'eau.
\end{enumerate}
}

\corr{11}{
\begin{enumerate}
  \item
  \[
  n(\ce{Cu})=\frac{12.7}{63.5}=\SI{0.200}{mol}.
  \]
  \item
  \[
  n(\ce{Zn})=\frac{6.50}{65.4}=\SI{9.94e-2}{mol}.
  \]
  \item
  \[
  x_{\ce{Cu}}=\frac{n(\ce{Cu})}{n(\ce{Cu})+n(\ce{Zn})}
  =\frac{0.200}{0.200+0.0994}=0.668.
  \]
\end{enumerate}
}

\corr{12}{
\begin{enumerate}
  \item
  \[
  M(\ce{C12H22O11})=12\times 12.0+22\times 1.0+11\times 16.0=\SI{342}{g.mol^{-1}}.
  \]
  \item
  \[
  C=\frac{c_m}{M}=\frac{34.2}{342}=\SI{0.100}{mol.L^{-1}}.
  \]
  \item
  \[
  V=\SI{0.2000}{L},\quad n=CV=0.100\times 0.2000=\SI{2.00e-2}{mol}.
  \]
\end{enumerate}
}

\corr{13}{
\begin{enumerate}
  \item Le rapport $A/C$ vaut toujours environ :
  \[
  \frac{0.12}{0.50}=0.24,\quad \frac{0.24}{1.00}=0.24,\quad \frac{0.36}{1.50}=0.24.
  \]
  Il y a proportionnalité.
  \item Avec $A=0.24C$ où $C$ est en \si{mmol.L^{-1}} :
  \[
  C=\frac{0.30}{0.24}=\SI{1.25}{mmol.L^{-1}}.
  \]
  \item Une limite possible est la non-linéarité de Beer-Lambert pour des solutions trop concentrées ou une absorbance trop élevée.
\end{enumerate}
}

\corr{14}{
\begin{enumerate}
  \item Oxydation du fer :
  \[
  \ce{Fe -> Fe^{2+} + 2e-}.
  \]
  \item Réduction du cuivre(II) :
  \[
  \ce{Cu^{2+} + 2e- -> Cu}.
  \]
  \item En additionnant :
  \[
  \ce{Fe + Cu^{2+} -> Fe^{2+} + Cu}.
  \]
  Les électrons se simplifient.
  \item L'oxydant est $\ce{Cu^{2+}}$ ; le réducteur est $\ce{Fe}$.
\end{enumerate}
}

\corr{15}{
\begin{enumerate}
  \item
  \[
  \begin{array}{c|ccc}
  & \ce{H2} & \ce{O2} & \ce{H2O}\\
  \hline
  n_i & 0.30 & 0.20 & 0\\
  \Delta n & -2x & -x & +2x\\
  n_f & 0.30-2x & 0.20-x & 2x
  \end{array}
  \]
  \item
  \[
  x\leq \frac{0.30}{2}=0.15,\qquad x\leq 0.20.
  \]
  Donc $x_{\max}=\SI{0.15}{mol}$.
  \item Le dihydrogène est limitant.
  \item
  \[
  n_f(\ce{H2})=0,\quad n_f(\ce{O2})=0.20-0.15=\SI{0.05}{mol},
  \]
  \[
  n_f(\ce{H2O})=2\times 0.15=\SI{0.30}{mol}.
  \]
\end{enumerate}
}

\corr{16}{
\begin{enumerate}
  \item Pour un mélange stœchiométrique :
  \[
  \frac{n(\ce{Al})}{2}=\frac{n(\ce{S})}{3}.
  \]
  \item
  \[
  n(\ce{S})=\frac{3}{2}n(\ce{Al})=\frac{3}{2}\times 0.40=\SI{0.60}{mol}.
  \]
  \item
  \[
  n(\ce{Al2S3})=\frac{n(\ce{Al})}{2}=\frac{0.40}{2}=\SI{0.20}{mol}.
  \]
\end{enumerate}
}

\corr{17}{
\begin{enumerate}
  \item À l'équivalence :
  \[
  n(\ce{H3O+})=n(\ce{HO-})_{\text{versé}}.
  \]
  \item
  \[
  n(\ce{HO-})=C_BV_E=0.100\times 12.5\times 10^{-3}=\SI{1.25e-3}{mol}.
  \]
  \item Cette quantité est celle d'acide dans \SI{20.0}{mL}. Donc :
  \[
  C_A=\frac{1.25\times 10^{-3}}{20.0\times 10^{-3}}=\SI{6.25e-2}{mol.L^{-1}}.
  \]
\end{enumerate}
}

\corr{18}{
\begin{enumerate}
  \item L'oxygène a 6 électrons de valence, l'hydrogène en a 1.
  \item Le schéma de Lewis de l'eau comporte deux liaisons O--H et deux doublets non liants sur l'oxygène.
  \item La géométrie est coudée.
  \item Les liaisons O--H sont polarisées vers O et la géométrie coudée ne permet pas la compensation des polarités : $\ce{H2O}$ est polaire.
\end{enumerate}
}

\corr{19}{
\begin{enumerate}
  \item
  \[
  C_{\text{apporté}}=\frac{n}{V}=\frac{0.050}{0.5000}=\SI{0.100}{mol.L^{-1}}.
  \]
  \item D'après l'équation, 1 mole de $\ce{CaCl2}$ donne 1 mole de $\ce{Ca^{2+}}$ :
  \[
  [\ce{Ca^{2+}}]=\SI{0.100}{mol.L^{-1}}.
  \]
  \item 1 mole de $\ce{CaCl2}$ donne 2 moles de $\ce{Cl-}$ :
  \[
  [\ce{Cl-}]=2C=\SI{0.200}{mol.L^{-1}}.
  \]
\end{enumerate}
}

\corr{20}{
\begin{enumerate}
  \item La chaîne principale comporte 3 atomes de carbone : préfixe \textit{propan-}.
  \item Le groupe caractéristique est le groupe hydroxyle $-\ce{OH}$ : famille des alcools.
  \item Le groupe $-\ce{OH}$ est porté par le carbone 1 : le nom est donc propan-1-ol.
\end{enumerate}
}

\corr{21}{
\begin{enumerate}
  \item Relation de dilution :
  \[
  C_0V_0=C_fV_f.
  \]
  \item Avec $V_f=\SI{100.0}{mL}$ :
  \[
  V_0=\frac{C_fV_f}{C_0}.
  \]
  Pour $\SI{1.0e-4}{mol.L^{-1}}$ :
  \[
  V_0=\frac{1.0\times 10^{-4}\times 100.0}{1.0\times 10^{-3}}=\SI{10.0}{mL}.
  \]
  Pour $\SI{2.0e-4}{mol.L^{-1}}$ :
  \[
  V_0=\SI{20.0}{mL}.
  \]
  Pour $\SI{4.0e-4}{mol.L^{-1}}$ :
  \[
  V_0=\SI{40.0}{mL}.
  \]
  \item Prélever \SI{20.0}{mL} de solution mère avec une pipette jaugée, les introduire dans une fiole jaugée de \SI{100.0}{mL}, compléter avec de l'eau distillée jusqu'au trait, boucher et homogénéiser.
\end{enumerate}
}

\corr{22}{
\begin{enumerate}
  \item La solution absorbe principalement dans le vert.
  \item Une solution qui absorbe le vert apparaît souvent de la couleur complémentaire, donc rouge-magenta.
  \item Au maximum d'absorption, la mesure est plus sensible : une petite variation de concentration produit une variation d'absorbance plus nette.
  \item
  \[
  C=\frac{A}{1250}=\frac{0.65}{1250}=\SI{5.2e-4}{mol.L^{-1}}.
  \]
\end{enumerate}
}

\corr{23}{
\begin{enumerate}
  \item
  \[
  \begin{array}{c|ccccc}
  & \ce{CaCO3} & \ce{H3O+} & \ce{Ca^{2+}} & \ce{CO2} & \ce{H2O}\\
  \hline
  n_i & 0.010 & 0.015 & 0 & 0 & \text{excès}\\
  \Delta n & -x & -2x & +x & +x & +3x\\
  n_f & 0.010-x & 0.015-2x & x & x & \text{excès}
  \end{array}
  \]
  \item Contraintes :
  \[
  x\leq 0.010,\qquad x\leq \frac{0.015}{2}=0.0075.
  \]
  Donc $x_{\max}=\SI{0.0075}{mol}$.
  \item Les ions $\ce{H3O+}$ sont limitants.
  \item
  \[
  n_f(\ce{CaCO3})=0.010-0.0075=\SI{0.0025}{mol}.
  \]
  \[
  n_f(\ce{H3O+})=0,\quad n_f(\ce{Ca^{2+}})=n_f(\ce{CO2})=\SI{0.0075}{mol}.
  \]
\end{enumerate}
}

\corr{24}{
\begin{enumerate}
  \item Puisqu'il reste \SI{0.20}{mol} de $B$ et que $n(B)=0.40-x$, on a :
  \[
  0.40-x_f=0.20,
  \]
  donc :
  \[
  x_f=\SI{0.20}{mol}.
  \]
  \item Pour une réaction $1:1$ :
  \[
  x_{\max}=\min(0.50,0.40)=\SI{0.40}{mol}.
  \]
  \item Comme $x_f<x_{\max}$, la transformation n'est pas totale.
  \item
  \[
  n_f(A)=0.50-0.20=\SI{0.30}{mol},
  \]
  \[
  n_f(C)=x_f=\SI{0.20}{mol}.
  \]
\end{enumerate}
}

\corr{25}{
\begin{enumerate}
  \item La demi-équation du fer doit être multipliée par 5 :
  \[
  \ce{5Fe^{2+} -> 5Fe^{3+} + 5e-}.
  \]
  \item En additionnant avec :
  \[
  \ce{MnO4- + 8H+ + 5e- -> Mn^{2+} + 4H2O},
  \]
  on obtient :
  \[
  \ce{MnO4- + 8H+ + 5Fe^{2+} -> Mn^{2+} + 4H2O + 5Fe^{3+}}.
  \]
  \item L'oxydant est $\ce{MnO4-}$ ; le réducteur est $\ce{Fe^{2+}}$.
\end{enumerate}
}

\corr{26}{
\begin{enumerate}
  \item À l'équivalence :
  \[
  n(\ce{Fe^{2+}})=5n(\ce{MnO4-}).
  \]
  \item
  \[
  n(\ce{MnO4-})=C_PV_E
  =2.00\times 10^{-2}\times 14.0\times 10^{-3}
  =\SI{2.80e-4}{mol}.
  \]
  \item
  \[
  n(\ce{Fe^{2+}})=5\times 2.80\times 10^{-4}=\SI{1.40e-3}{mol}.
  \]
  \item
  \[
  C(\ce{Fe^{2+}})=\frac{1.40\times 10^{-3}}{10.0\times 10^{-3}}
  =\SI{0.140}{mol.L^{-1}}.
  \]
\end{enumerate}
}

\corr{27}{
\begin{enumerate}
  \item Le carbone est central, lié par deux doubles liaisons aux deux oxygènes :
  \[
  \ce{O=C=O}.
  \]
  Chaque oxygène possède deux doublets non liants.
  \item La molécule est linéaire.
  \item Les liaisons C=O sont polarisées vers l'oxygène, plus électronégatif.
  \item Les deux polarités de liaison sont opposées et de même intensité dans une géométrie linéaire symétrique. Elles se compensent : $\ce{CO2}$ est apolaire.
\end{enumerate}
}

\corr{28}{
\begin{enumerate}
  \item L'azote est lié à trois hydrogènes et porte un doublet non liant.
  \item La géométrie est pyramidale trigonale.
  \item Les liaisons N--H sont polarisées vers N et la géométrie pyramidale empêche la compensation : $\ce{NH3}$ est polaire.
  \item Dans $\ce{BF3}$, les trois polarités se compensent par symétrie plane triangulaire ; la molécule est donc apolaire.
\end{enumerate}
}

\corr{29}{
\begin{enumerate}
  \item On utilise le cyclohexane.
  \item L'ampoule à décanter permet d'agiter les deux phases puis de les séparer après décantation.
  \item Le diiode se trouve majoritairement dans la phase organique, le cyclohexane.
  \item Le diiode est apolaire et se dissout mieux dans un solvant apolaire : c'est l'idée « qui se ressemble s'assemble ».
\end{enumerate}
}

\corr{30}{
\begin{enumerate}
  \item On ajuste :
  \[
  \ce{C2H6O + 3O2 -> 2CO2 + 3H2O}.
  \]
  \item D'après l'équation, 1 mole d'éthanol forme 2 moles de $\ce{CO2}$ :
  \[
  n(\ce{CO2})=2\times 0.50=\SI{1.00}{mol}.
  \]
  \item 1 mole d'éthanol consomme 3 moles de $\ce{O2}$ :
  \[
  n(\ce{O2})=3\times 0.50=\SI{1.50}{mol}.
  \]
  \item Le dioxyde de carbone est un gaz à effet de serre ; la combustion contribue donc aux émissions de $\ce{CO2}$.
\end{enumerate}
}

\corr{31}{
\begin{enumerate}
  \item
  \[
  n(\ce{I2})=C_IV_E=5.00\times 10^{-3}\times 18.2\times 10^{-3}
  =\SI{9.10e-5}{mol}.
  \]
  \item La réaction est de proportion $1:1$, donc :
  \[
  n_{\text{prélèvement}}(\ce{C6H8O6})=\SI{9.10e-5}{mol}.
  \]
  \item La fiole de \SI{250.0}{mL} contient $250.0/20.0=12.5$ fois plus :
  \[
  n_{\text{fiole}}=12.5\times 9.10\times 10^{-5}
  =\SI{1.14e-3}{mol}.
  \]
  \item
  \[
  m=nM=1.14\times 10^{-3}\times 176.0=\SI{0.200}{g}=\SI{200}{mg}.
  \]
  \item La valeur mesurée est très inférieure à \SI{500}{mg}. Le comprimé analysé ne correspond pas à cette indication ou l'hypothèse/protocole comporte une erreur, dilution, dégradation ou perte.
\end{enumerate}
}

\corr{32}{
\begin{enumerate}
  \item Les ions permanganate $\ce{MnO4-}$ sont violets. Avant l'équivalence, ils sont consommés ; après l'équivalence, un léger excès colore durablement la solution.
  \item
  \[
  n(\ce{MnO4-})=CV_E=1.00\times 10^{-3}\times 9.60\times 10^{-3}
  =\SI{9.60e-6}{mol}.
  \]
  \item À l'équivalence :
  \[
  n(\ce{Fe^{2+}})=5n(\ce{MnO4-})=\SI{4.80e-5}{mol}.
  \]
  Dans \SI{50.0}{mL} :
  \[
  C(\ce{Fe^{2+}})=\frac{4.80\times 10^{-5}}{50.0\times 10^{-3}}
  =\SI{9.60e-4}{mol.L^{-1}}.
  \]
  \item
  \[
  c_m=CM=9.60\times 10^{-4}\times 55.8
  =\SI{5.36e-2}{g.L^{-1}}.
  \]
\end{enumerate}
}

\corr{33}{
\begin{enumerate}
  \item
  \[
  n(A)=0.80-2x,\quad n(B)=0.90-3x,\quad n(C)=x,\quad n(D)=4x.
  \]
  \item
  \[
  x\leq \frac{0.80}{2}=0.40,\qquad x\leq \frac{0.90}{3}=0.30.
  \]
  Donc :
  \[
  x_{\max}=\SI{0.30}{mol}.
  \]
  \item
  \[
  n_f(A)=0.80-2\times 0.30=\SI{0.20}{mol},
  \]
  \[
  n_f(B)=0,\quad n_f(C)=\SI{0.30}{mol},\quad n_f(D)=\SI{1.20}{mol}.
  \]
  \item Pseudo-code :
\begin{verbatim}
xA = nA0/a
xB = nB0/b
si xA < xB alors A est limitant et xmax = xA
sinon B est limitant et xmax = xB
calculer les quantités finales
\end{verbatim}
\end{enumerate}
}

\corr{34}{
\begin{enumerate}
  \item L'absorption est maximale autour de \SI{550}{nm}.
  \item \SI{550}{nm} correspond au vert-jaune.
  \item La couleur transmise est approximativement complémentaire : la solution peut apparaître rouge-violette selon le spectre transmis.
  \item À $\lambda=\SI{550}{nm}$, l'absorbance est maximale, donc la mesure est plus sensible.
\end{enumerate}
}

\corr{35}{
\begin{enumerate}
  \item Le bore a 3 électrons de valence, l'hydrogène en a 1.
  \item Le bore forme trois liaisons simples B--H.
  \item Le bore est entouré de seulement 6 électrons dans ce schéma, pas 8.
  \item Une lacune électronique correspond à un déficit d'électrons autour d'un atome, qui peut accepter un doublet.
  \item La molécule plane triangulaire est symétrique ; les polarités de liaison se compensent : elle est apolaire.
\end{enumerate}
}

\corr{36}{
\begin{enumerate}
  \item
  \[
  \ce{Al2(SO4)3(s) -> 2Al^{3+}(aq) + 3SO4^{2-}(aq)}.
  \]
  \item
  \[
  M=2\times 27.0+3(32.1+4\times 16.0)=\SI{342.3}{g.mol^{-1}}.
  \]
  \item
  \[
  n=\frac{3.42}{342.3}=\SI{9.99e-3}{mol}.
  \]
  \[
  C=\frac{n}{V}=\frac{9.99\times 10^{-3}}{0.2500}=\SI{4.00e-2}{mol.L^{-1}}.
  \]
  \item
  \[
  [\ce{Al^{3+}}]=2C=\SI{8.00e-2}{mol.L^{-1}},
  \]
  \[
  [\ce{SO4^{2-}}]=3C=\SI{1.20e-1}{mol.L^{-1}}.
  \]
\end{enumerate}
}

\corr{37}{
\begin{enumerate}
  \item La longue chaîne carbonée $\ce{CH3-(CH2)16-}$ est hydrophobe.
  \item Le groupe carboxylate $\ce{-COO-}$ est hydrophile.
  \item L'entité possède une partie compatible avec l'eau et une partie compatible avec les graisses : elle est amphiphile.
  \item Les queues hydrophobes s'associent à la graisse tandis que les têtes hydrophiles restent au contact de l'eau. Cela favorise la dispersion des graisses sous forme de micelles.
\end{enumerate}
}

\corr{38}{
\begin{enumerate}
  \item La forte bande vers \SI{1715}{cm^{-1}} met en évidence une liaison carbonyle C=O.
  \item L'absence de large bande O--H permet d'exclure un alcool.
  \item Les familles compatibles sont aldéhyde ou cétone.
  \item Une formule possible est $\ce{CH3COCH3}$, la propanone.
\end{enumerate}
}

\corr{39}{
\begin{enumerate}
  \item La réaction étant $1:1$, l'acide carboxylique, introduit à \SI{0.20}{mol}, est limitant.
  \item
  \[
  n_{\max}(\text{ester})=\SI{0.20}{mol}.
  \]
  \item
  \[
  n_{\text{obt}}=\frac{17.6}{116}=\SI{0.152}{mol}.
  \]
  \item
  \[
  \eta=\frac{0.152}{0.20}\times 100=76.0\%.
  \]
\end{enumerate}
}

\corr{40}{
\begin{enumerate}
  \item Liaisons rompues : 4 liaisons C--H dans $\ce{CH4}$ et 2 liaisons O=O dans $\ce{O2}$.
  \[
  E_{\text{rompues}}=4\times 413+2\times 498=\SI{2648}{kJ.mol^{-1}}.
  \]
  \item Liaisons formées : 2 liaisons C=O dans $\ce{CO2}$ et 4 liaisons O--H dans $2\ce{H2O}$.
  \[
  E_{\text{formées}}=2\times 799+4\times 463=\SI{3450}{kJ.mol^{-1}}.
  \]
  \item
  \[
  \Delta_rE=E_{\text{rompues}}-E_{\text{formées}}
  =2648-3450=\SI{-802}{kJ.mol^{-1}}.
  \]
  \item L'énergie est négative : la réaction libère de l'énergie, elle est exothermique.
\end{enumerate}
}

\corr{41}{
\begin{enumerate}
  \item
  \[
  n_0(\ce{H3O+})=CV=1.00\times 0.1000=\SI{0.100}{mol}.
  \]
  \item Il reste \SI{0.020}{mol}, donc la quantité consommée vaut :
  \[
  n_{\text{cons}}(\ce{H3O+})=0.100-0.020=\SI{0.080}{mol}.
  \]
  Comme $\Delta n(\ce{H3O+})=-2x$ :
  \[
  2x_f=0.080 \Rightarrow x_f=\SI{0.040}{mol}.
  \]
  \item Le tartre a totalement disparu, donc :
  \[
  n_0(\ce{CaCO3})=x_f=\SI{0.040}{mol}.
  \]
  \item
  \[
  m= nM=0.040\times 100.1=\SI{4.00}{g}.
  \]
  \item
  \[
  n(\ce{CO2})=x_f=\SI{0.040}{mol}.
  \]
  \[
  V=nV_m=0.040\times 24.0=\SI{0.96}{L}.
  \]
  \item Le dégagement gazeux est dû à la formation de dioxyde de carbone.
\end{enumerate}
}

\corr{42}{
\begin{enumerate}
  \item Les points sont presque alignés et l'absorbance augmente avec la concentration : la loi de Beer-Lambert est utilisable dans cette gamme.
  \item On peut estimer :
  \[
  k\simeq \frac{0.76}{10.0}=\SI{0.076}{L.mg^{-1}}.
  \]
  \item
  \[
  C_{\text{diluée}}=\frac{A}{k}=\frac{0.54}{0.076}=\SI{7.1}{mg.L^{-1}}.
  \]
  \item La boisson initiale est 5 fois plus concentrée :
  \[
  C_0=5\times 7.1=\SI{35.5}{mg.L^{-1}}.
  \]
  \item Une dilution permet de ramener l'absorbance dans le domaine de validité de la loi de Beer-Lambert et d'éviter une solution trop absorbante.
\end{enumerate}
}

\corr{43}{
\begin{enumerate}
  \item L'oxydant titrant est l'ion permanganate $\ce{MnO4-}$.
  \item À l'équivalence :
  \[
  \frac{n(\ce{MnO4-})}{2}=\frac{n(\ce{H2O2})}{5}.
  \]
  \item
  \[
  n(\ce{MnO4-})=CV_E=2.00\times 10^{-2}\times 16.0\times 10^{-3}
  =\SI{3.20e-4}{mol}.
  \]
  \item
  \[
  n(\ce{H2O2})=\frac{5}{2}n(\ce{MnO4-})
  =\frac{5}{2}\times 3.20\times 10^{-4}
  =\SI{8.00e-4}{mol}.
  \]
  \item
  \[
  C(\ce{H2O2})=\frac{8.00\times 10^{-4}}{10.0\times 10^{-3}}
  =\SI{8.00e-2}{mol.L^{-1}}.
  \]
  \item Si la solution commerciale a été diluée 10 fois :
  \[
  C_{\text{commerciale}}=10C=\SI{0.800}{mol.L^{-1}}.
  \]
\end{enumerate}
}

\corr{44}{
\begin{enumerate}
  \item On multiplie l'oxydation du fer par 2 :
  \[
  \ce{2Fe -> 2Fe^{2+} + 4e-}.
  \]
  En additionnant :
  \[
  \ce{2Fe + O2 + 4H+ -> 2Fe^{2+} + 2H2O}.
  \]
  \item L'oxydant est $\ce{O2}$ ; le réducteur est $\ce{Fe}$.
  \item
  \[
  n(\ce{Fe})=\frac{1.12}{55.8}=\SI{2.01e-2}{mol}.
  \]
  \item D'après l'équation, $2$ moles de Fe consomment $1$ mole de $\ce{O2}$ :
  \[
  n(\ce{O2})=\frac{n(\ce{Fe})}{2}=\SI{1.00e-2}{mol}.
  \]
  \item
  \[
  V(\ce{O2})=nV_m=1.00\times 10^{-2}\times 24.0=\SI{0.240}{L}.
  \]
  \item L'eau permet les échanges ioniques et la présence d'espèces dissoutes ; elle favorise donc les processus électrochimiques de corrosion.
\end{enumerate}
}

\corr{45}{
\begin{enumerate}
  \item On choisit l'éthanoate d'éthyle car $S$ y est beaucoup plus soluble que dans l'eau.
  \item Si les deux liquides étaient miscibles, il n'y aurait pas deux phases séparables.
  \item $S$ se trouve majoritairement dans la phase organique.
  \item Plusieurs petites extractions sont souvent plus efficaces car à chaque extraction une fraction de soluté passe dans le solvant organique frais.
  \item Introduire la solution aqueuse et le solvant organique dans l'ampoule, agiter en dégazant, laisser décanter, identifier les phases, récupérer la phase contenant le soluté.
\end{enumerate}
}

\corr{46}{
\begin{enumerate}
  \item L'eau est coudée et possède deux liaisons O--H polarisées vers O. Les polarités ne se compensent pas.
  \item Dans l'éthanol, le groupe $-\ce{OH}$ est polaire et hydrophile ; la chaîne $\ce{CH3CH2-}$ est plutôt apolaire.
  \item L'éthanol forme des interactions favorables avec l'eau, notamment des ponts hydrogène.
  \item L'hexane est apolaire et ne forme pas d'interactions favorables avec l'eau polaire.
  \item Une espèce très apolaire sera plus soluble dans l'hexane.
  \item Les espèces présentant des polarités et interactions semblables se mélangent plus facilement.
\end{enumerate}
}

\corr{47}{
\begin{enumerate}
  \item Transformation : chauffage à reflux. Isolement : séparation du produit du milieu réactionnel. Purification : lavage, séchage, distillation ou recristallisation. Analyse : CCM, température de fusion, spectroscopie.
  \item La réaction est $1:1$ ; l'acide, introduit à \SI{0.150}{mol}, est limitant.
  \item
  \[
  n_{\max}(\text{ester})=\SI{0.150}{mol}.
  \]
  \item
  \[
  n_{\text{obt}}=\frac{12.8}{102}=\SI{0.125}{mol}.
  \]
  \item
  \[
  \eta=\frac{0.125}{0.150}\times 100=\SI{83.3}{\percent}.
  \]
  \item Pertes lors des transferts, réaction non totale, purification imparfaite, produit resté dans la phase aqueuse ou évaporation partielle.
\end{enumerate}
}

\corr{48}{
\begin{enumerate}
  \item La bande large entre \SI{3200}{cm^{-1}} et \SI{3600}{cm^{-1}} met en évidence une liaison O--H.
  \item La famille probable est celle des alcools.
  \item Deux formules possibles :
  \[
  \ce{CH3-CH2-CH2-OH}
  \]
  et
  \[
  \ce{CH3-CH(OH)-CH3}.
  \]
  \item Les noms sont propan-1-ol et propan-2-ol.
  \item Oui, elles ont la même formule brute $\ce{C3H8O}$.
  \item Ce sont des isomères.
\end{enumerate}
}

\corr{49}{
\begin{enumerate}
  \item
  \[
  E=m c \Delta T=500\times 4.18\times 18.0
  =\SI{3.76e4}{J}=\SI{37.6}{kJ}.
  \]
  \item
  \[
  n=\frac{2.30}{46.0}=\SI{5.00e-2}{mol}.
  \]
  \item
  \[
  E_m=\frac{37.6}{5.00\times 10^{-2}}=\SI{752}{kJ.mol^{-1}}.
  \]
  \item
  \[
  PC=\frac{37.6}{2.30}=\SI{16.3}{kJ.g^{-1}}.
  \]
  \item Une partie de l'énergie chauffe l'air, le récipient ou est perdue par rayonnement ; la combustion peut aussi être incomplète.
\end{enumerate}
}

\corr{50}{
\begin{enumerate}
  \item L'oxydant est $\ce{Cu^{2+}}$ et le réducteur est $\ce{Zn}$.
  \item
  \[
  \ce{Zn -> Zn^{2+} + 2e-},
  \qquad
  \ce{Cu^{2+} + 2e- -> Cu}.
  \]
  \item
  \[
  n(\ce{Zn})=\frac{1.31}{65.4}=\SI{2.00e-2}{mol}.
  \]
  \[
  n(\ce{Cu^{2+}})=CV=0.150\times 0.1000=\SI{1.50e-2}{mol}.
  \]
  \item La réaction est $1:1$ ; les ions cuivre(II), moins nombreux, sont limitants.
  \item
  \[
  n(\ce{Cu})=1.50\times 10^{-2}\ \text{mol}.
  \]
  \[
  m(\ce{Cu})=nM=1.50\times 10^{-2}\times 63.5=\SI{0.953}{g}.
  \]
  \item Le zinc consommé vaut \SI{1.50e-2}{mol}. Il reste :
  \[
  n_f(\ce{Zn})=2.00\times 10^{-2}-1.50\times 10^{-2}
  =\SI{5.00e-3}{mol}.
  \]
  \[
  m_f(\ce{Zn})=5.00\times 10^{-3}\times 65.4=\SI{0.327}{g}.
  \]
  \item Dans une pile, les électrons libérés par l'oxydation du zinc circulent dans le circuit extérieur vers l'espèce oxydante. Ce déplacement ordonné d'électrons constitue un courant électrique.
\end{enumerate}
}

\corr{Bilan 1}{
\begin{enumerate}
  \item
  \[
  C_{\text{diluée}}=\frac{A}{120}=\frac{0.36}{120}=\SI{3.00e-3}{mol.L^{-1}}.
  \]
  \item La solution initiale est 4 fois plus concentrée :
  \[
  C_0=4C_{\text{diluée}}=\SI{1.20e-2}{mol.L^{-1}}.
  \]
  \item
  \[
  n(\ce{Cu^{2+}})=C_0V=1.20\times 10^{-2}\times 0.2500
  =\SI{3.00e-3}{mol}.
  \]
  \item La dissolution de $\ce{CuSO4}$ donne une mole de $\ce{Cu^{2+}}$ par mole de soluté :
  \[
  n(\ce{CuSO4})=\SI{3.00e-3}{mol}.
  \]
  \[
  m=nM=3.00\times 10^{-3}\times 159.6=\SI{0.479}{g}.
  \]
\end{enumerate}
}

\corr{Bilan 2}{
\begin{enumerate}
  \item
  \[
  n(\ce{HO-})=C_BV_E=0.200\times 25.0\times 10^{-3}
  =\SI{5.00e-3}{mol}.
  \]
  \item La réaction est $1:1$, donc :
  \[
  n(\ce{AH})=\SI{5.00e-3}{mol}.
  \]
  \item
  \[
  M=\frac{m}{n}=\frac{0.620}{5.00\times 10^{-3}}
  =\SI{124}{g.mol^{-1}}.
  \]
  \item
  \[
  \eta=\frac{n_{\text{exp}}}{n_{\max}}\times 100
  =\frac{5.00\times 10^{-3}}{6.20\times 10^{-3}}\times 100
  =80.6\%.
  \]
  \item On peut vérifier la pureté par chromatographie sur couche mince, mesure de température de fusion, spectroscopie infrarouge ou comparaison avec une référence.
\end{enumerate}
}

\newpage

%====================================================
\section{Fiche méthode finale}

\begin{retenirbox}
\textbf{Grandeurs de composition}
\[
n=\frac{m}{M},\qquad n=\frac{V}{V_m},\qquad C=\frac{n}{V},\qquad c_m=CM.
\]
Toujours convertir les volumes en litres lorsque la concentration est exprimée en \si{mol.L^{-1}}.
\end{retenirbox}

\begin{methodbox}
\textbf{Tableau d'avancement}
\begin{enumerate}
  \item Ajuster l'équation de réaction.
  \item Écrire les quantités initiales.
  \item Introduire l'avancement $x$.
  \item Écrire les quantités finales en fonction de $x$.
  \item Déterminer $x_{\max}$ avec les réactifs.
  \item Si la transformation est totale, $x_f=x_{\max}$.
  \item Si elle n'est pas totale, $x_f<x_{\max}$.
\end{enumerate}
\end{methodbox}

\begin{retenirbox}
\textbf{Oxydo-réduction}
Un oxydant capte des électrons ; un réducteur cède des électrons.
Pour établir une réaction redox :
\begin{enumerate}
  \item écrire les deux demi-équations ;
  \item équilibrer les électrons ;
  \item additionner ;
  \item simplifier les électrons.
\end{enumerate}
\end{retenirbox}

\begin{methodbox}
\textbf{Titrage}
À l'équivalence, les réactifs sont introduits en proportions stœchiométriques.
Pour :
\[
aA+bB\to \text{produits},
\]
on a :
\[
\frac{n(A)}{a}=\frac{n_E(B)}{b}.
\]
\end{methodbox}

\begin{retenirbox}
\textbf{Lewis, géométrie et polarité}
Une molécule peut avoir des liaisons polarisées sans être polaire si les polarités se compensent par symétrie. Exemple : $\ce{CO2}$ est apolaire bien que les liaisons C=O soient polarisées.
\end{retenirbox}

\begin{retenirbox}
\textbf{Dissolution d'un solide ionique}
\[
\ce{NaCl(s) -> Na+(aq) + Cl-(aq)}
\]
\[
\ce{CaCl2(s) -> Ca^{2+}(aq) + 2Cl-(aq)}
\]
Les coefficients de l'équation de dissolution permettent de déterminer les concentrations ioniques.
\end{retenirbox}

\begin{retenirbox}
\textbf{Chimie organique}
\begin{itemize}
  \item alcool : groupe hydroxyle $-\ce{OH}$ ;
  \item aldéhyde : groupe carbonyle en bout de chaîne $-\ce{CHO}$ ;
  \item cétone : groupe carbonyle au milieu de chaîne ;
  \item acide carboxylique : groupe $-\ce{COOH}$.
\end{itemize}
\end{retenirbox}

\begin{methodbox}
\textbf{Rendement}
\[
\eta=\frac{n_{\text{obtenu}}}{n_{\max}}\times 100.
\]
Un rendement inférieur à 100\% peut venir d'une réaction non totale, de pertes lors des manipulations, d'une purification imparfaite ou d'une dégradation du produit.
\end{methodbox}

\begin{retenirbox}
\textbf{Combustion et énergie}
Une combustion complète d'une espèce organique produit $\ce{CO2}$ et $\ce{H2O}$.
L'énergie de réaction peut être estimée par :
\[
\Delta_rE\simeq E_{\text{liaisons rompues}}-E_{\text{liaisons formées}}.
\]
Si $\Delta_rE<0$, la transformation est exothermique.
\end{retenirbox}

\begin{erreurbox}
\textbf{Erreurs fréquentes à éviter}
\begin{itemize}
  \item oublier les coefficients stœchiométriques dans un titrage ;
  \item confondre concentration molaire et concentration massique ;
  \item utiliser des millilitres dans $n=CV$ sans convertir ;
  \item oublier qu'un ion solide dissous peut donner plusieurs ions en solution ;
  \item conclure qu'une molécule est polaire seulement parce qu'elle possède des liaisons polarisées ;
  \item calculer un rendement avec une masse au lieu d'une quantité de matière lorsque les masses molaires diffèrent ;
  \item oublier d'ajuster une équation de combustion avant tout calcul.
\end{itemize}
\end{erreurbox}

\begin{center}
\Large\bfseries Fin du TD
\end{center}

\end{document}
```
